\documentclass[11pt]{article}
\usepackage{amsmath,amssymb}
\usepackage{lmodern}
\usepackage[T1]{fontenc}
\usepackage[margin=1.1in]{geometry}
\newcommand{\kp}{k_\perp}
\newcommand{\gpar}{\gamma_\parallel}
\newcommand{\nuin}{\nu_{in}}
\title{Derivation of the GDI linear matrix $L$ in \texttt{jax\_rmhd}\\
\large and resolution of the apparent discrepancy with eq.~(5.3) of the notes}
\author{P4a implementation note}
\date{2026-08-01}
\begin{document}
\maketitle

\noindent All equation numbers refer to \texttt{docs/GDI\_nonlinear\_equations (10).pdf}.
The code evolves $f=(N,\varphi)$ with the convention
$\partial_t f = L f + \mathcal{N}(f)$, propagator $e^{L\tau}$ (putzer2 backend).
Everything below is in the normalized units of \S5.1: $\rho_s=c_s=\Omega_i=1$.

\section{Starting point}

The normalized 2D equations (5.4)--(5.5), with the parallel current-closure term
represented by a perpendicular-$k$-local $\gpar$:
\begin{align}
\frac{\partial N}{\partial t} - \frac{1}{L_n}\frac{\partial\varphi}{\partial y}
  + \{\varphi, N\} - \gpar\,[\varphi - N] &= 0,
  \tag{5.4}\\
\frac{\partial \nabla_\perp^2\varphi}{\partial t} + \{\varphi,\nabla_\perp^2\varphi\}
  + \nuin\nabla_\perp^2\varphi - \nuin v_0 \frac{\partial N}{\partial y}
  - \gpar\,[\varphi-N] &= 0.
  \tag{5.5}
\end{align}
For P4a (strictly 2D) the closure floor of eq.~(4.3) is used, scaled by the tunable
\texttt{eqpars} key \texttt{gpar\_fac}:
\begin{equation}
\gpar = \texttt{gpar\_fac}\;\nuin \kp^2 \;(\text{default } 1).
\label{eq:closure}
\end{equation}
(Relation to eq.~(4.7)'s $\alpha$: with $\gpar$ substituted for $D_\parallel k_z^2$ one
gets $\alpha = 1 + \texttt{gpar\_fac}$, so the default reproduces $\alpha=2$, the
minimum-$\gpar$ factor quoted at eq.~(4.6). Below, $\alpha$ denotes
$\texttt{gpar\_fac}$ itself only inside the matrix \eqref{eq:L}, where it abbreviates
$\gpar/\kp^2\nuin$.)
Note $\gpar/\kp^2 = \alpha\nuin$ is $k$-independent under this closure --- this constant
appears repeatedly below.

\section{Fourier form and the vorticity division}

With $\partial_y \to i k_y$ and $\nabla_\perp^2 \to -\kp^2$, the linear part of (5.4) is
immediate:
\begin{equation}
\partial_t N = \underbrace{-\gpar}_{L_{NN}} N
 + \underbrace{\Big(\gpar + \tfrac{i k_y}{L_n}\Big)}_{L_{N\varphi}} \varphi
 \;+\; \text{(bracket)}.
\end{equation}
Equation (5.5) evolves the vorticity $-\kp^2\varphi$; dividing by $-\kp^2$
(zero mode masked via \texttt{inv\_ksq}, exactly as the RMHD $\varphi$ equation):
\begin{equation}
\partial_t \varphi
 = \underbrace{\Big(\frac{\gpar}{\kp^2} - \frac{i k_y \nuin v_0}{\kp^2}\Big)}_{L_{\varphi N}} N
 + \underbrace{\Big(-\nuin - \frac{\gpar}{\kp^2}\Big)}_{L_{\varphi\varphi}} \varphi
 \;+\; \frac{\{\varphi,\nabla_\perp^2\varphi\}}{-\kp^2}\Big|_{\text{sign as coded}}.
\end{equation}
Collecting, and adding the numerical hyperdissipation closure on both diagonal entries
(required because the linear growth rate does not decay at large $\kp$, cf.\ eq.~3.14):
\begin{equation}
\boxed{\;
L \;=\;
\begin{pmatrix}
 -\gpar - D_k & \gpar + \dfrac{i k_y}{L_n}\\[2ex]
 \alpha\nuin - \dfrac{i k_y \nuin v_0}{\kp^2} & -\nuin - \alpha\nuin - D_k
\end{pmatrix},
\qquad D_k \equiv \texttt{diss}\,\kp^{2\cdot\texttt{hyper}},
\;}
\label{eq:L}
\end{equation}
using $\gpar/\kp^2=\alpha\nuin$ from \eqref{eq:closure}. This is verbatim
\texttt{\_L\_entries} in \texttt{jax\_rmhd/physics/gdi.py}. (The $ik_y$ factors use
$k_y$ with the Nyquist row zeroed: an odd-order spectral derivative is ill-defined on the
self-conjugate row, and unlike the nonlinear terms, $L$ is applied without dealiasing, so
the row must be fixed explicitly to preserve rfft2 reality. One edge row, no interior
physics.)

\section{Check against the dispersion relation}

With $f\propto e^{\lambda t}$ and $D_k=0$, the characteristic polynomial of \eqref{eq:L},
$\lambda^2 - (\operatorname{tr}L)\lambda + \det L = 0$, factors as
\begin{equation}
(\lambda + \gpar)\Big[(\lambda + \nuin)\kp^2 + \gpar\Big]
 = \big(\gpar + i\omega_*\big)\big(\gpar - i\omega_{\rm Ped}\big),
\qquad
\omega_* \equiv \frac{k_y}{L_n},\;\;
\omega_{\rm Ped} \equiv k_y \nuin v_0 .
\end{equation}
Substituting $\lambda = -i\omega$ (each bracket picks up a factor $-i$; the right side
rearranges with $i\cdot(-i)=1$) gives exactly eq.~(3.7) of the notes:
\begin{equation}
(\omega + i\gpar)\Big[(\omega + i\nuin)\kp^2 + i\gpar\Big]
 = -\big(\omega_{\rm Ped} + i\gpar\big)\big(\omega_* - i\gpar\big).
\tag{3.7}
\end{equation}
At $\gpar=0$ this collapses to $\omega(\omega+i\nuin) = -\omega_*\omega_{\rm Ped}/\kp^2
= -A$ with $A = k_y^2\nuin v_0/(\kp^2 L_n)$, i.e.\ eq.~(2.8); the code's tests further
verify the collisional limit (2.9), the inertial limit (2.11), and, at finite $\gpar$,
the roots of (3.7) directly (numpy \texttt{roots} vs.\ eigenvalues of \eqref{eq:L},
rtol $10^{-10}$), plus a propagator-level check that putzer2-evolved fields grow at
$\max\operatorname{Re}(m\pm s)$.

\section{The ``corrected term'' that wasn't}

The P4a implementation report claimed the derived $L_{\varphi N}$ contains an
$\alpha\nuin$ term \emph{absent} from the notes' transcribed matrix (5.3). That claim is
wrong, and the note-writer (not the paper) needed correcting: on visual inspection of
p.~7, eq.~(5.3) reads, in the $\partial_t f + M f = 0$ convention,
\begin{equation}
M = \begin{pmatrix}
\gpar & -i k_y \rho_s c_s/L_n - \gpar\\[0.5ex]
i k_y \nuin v_0 / \kp^2\rho_s c_s \;\underline{-\;\gpar/\kp^2\rho_s^2} &
\nuin + \gpar/\kp^2\rho_s^2
\end{pmatrix},
\tag{5.3}
\end{equation}
whose $(\varphi,N)$ entry \emph{does} include $-\gpar/\kp^2\rho_s^2$. Hence
$L = -M$ agrees with \eqref{eq:L} in \emph{every} entry (before the added
hyperdissipation): the derivation from (5.4)--(5.5), the transcription (5.3), and the
implemented matrix are all mutually consistent. The implementation's validation never
relied on (5.3) --- the dispersion-relation tests above are the arbiter --- but no
discrepancy with the notes exists.

\section{Sanity limits}

$\gpar\to\infty$ (adiabatic): the $N$-equation row forces $N\to\varphi$
(eq.~3.8 with $\gpar$ dominant), and the growth shuts off --- consistent with the
stabilization role of current closure. $\nuin$ enters $L_{\varphi\varphi}$ as bare
vorticity friction with relaxation rate exactly $\nuin$ (the ``mild'' stiffness claim in
the plan's design background). Growing solutions at $A>0$: $s^2 = m^2 - \det L$ is
complex in general; all propagator arithmetic is complex, and the growth rate extracted
from evolved fields matches $\operatorname{Re}(m+s)$ to $10^{-6}$ (fp64).

\end{document}
